% tikz-tensors-nodes.tex -- node types: what a single tensor looks like, and
% what the layout and the contraction rules need to know about it.
%
% A node type is a TikZ style built on `tn node'. It says three things:
%   its outline   a TikZ shape, a fill (paper unless a convention fills it) and
%                 the size at which it reads as the same weight as the others
%   its size      tn wide, tn tall, tn small, tn flat resize any outline,
%                 and tn size=<factor> scales it
%   its indices   where they leave it, as properties a layout reads back
%                 when it places the node:
%     tn site anchor=<anchor> the point its index down the site leaves from,
%                             and the point a stack puts on the site's line
%                             (center unless the shape says otherwise)
%     tn flux={<side>=in|out, ...}
%                             which way its indices on a side run: into the
%                             tensor or out of it. A bond runs from the end
%                             that is out to the end that is in -- the way an
%                             arrow on it points (tn arrow) -- and two ends
%                             both in, or both out, are an error. A triangle
%                             leans one way by itself -- in at its flat side,
%                             out at its apex -- which the other end of a bond,
%                             or a type's own tn flux, may overrule.
%     tn single=up|down       the side on which it has one index however many
%                             sites it spans (the apex of a triangle)
%     tn index={<name>:<side>[:in|out], ...}
%                             its indices, each named and leaving on a side --
%                             up, down, left, right (on a grid: nw, ne, se,
%                             sw, down) -- in an order, which is how a figure
%                             may count them: T:1 is the first. Several may
%                             leave on one side, {a:left, b:left}. Not said, a
%                             tensor has the indices its layout gives it, one
%                             on each side per site or layer it covers; said,
%                             it has these and no others, and a picture that
%                             ends with one of them not drawn is an error.
%                             An index may say which way it runs, which comes
%                             before its side's (tn flux).
% A user's own type is a style of the same kind and is a type like these: a
% kite pointing up is {tn node, kite, tn fill=green3, tn single=up}.
% The package's types are named for their shape and never for what they mean;
% a meaning is a style of the user's notation built on one of them.
% tn size=<f> draws an outline at <f> times the size its type gives it -- a
% factor and not a length, so a figure may say it: a tensor of a compact
% picture, or one that is to stand out. The label keeps its size, which is the
% size of the equations beside it. A layout says it for a block or a layer
% drawn smaller (size=, scale=; tikz-tensors-layout.tex).
%
% An outline with rounded corners says its radius with tn@radius, so that
% tn size scales the radius with the outline: a capsule drawn smaller stays a
% capsule instead of a box whose corners overrun its sides.
\tikzset{tn@radius/.code={\def\tn@rad{#1}\tikzset{rounded corners=#1}},
  tn size/.code={%
  \pgfmathsetlengthmacro\tn@szw{(\pgfkeysvalueof{/pgf/minimum width})*(#1)}%
  \pgfmathsetlengthmacro\tn@szh{(\pgfkeysvalueof{/pgf/minimum height})*(#1)}%
  \pgfkeyslet{/pgf/minimum width}\tn@szw \pgfkeyslet{/pgf/minimum height}\tn@szh
  \ifx\tn@rad\@empty\else
    \pgfmathsetlengthmacro\tn@szr{(\tn@rad)*(#1)}\tikzset{rounded corners=\tn@szr}\fi}}
\def\tn@rad{}
\gdef\tn@legat{center}
\gdef\tn@flux{}
\gdef\tn@fluxw{}
\gdef\tn@ixdirs{}
\gdef\tn@single{}
\gdef\tn@legs{}
\gdef\tn@rank{}
\gdef\tn@index{}
\gdef\tn@tname{}
\tikzset{tn site anchor/.code={\gdef\tn@legat{#1}},
         tn flux/.code={\xdef\tn@flux{\zap@space#1 \@empty}},
         tn@lean/.code={\xdef\tn@fluxw{\zap@space#1 \@empty}},
         tn single/.code={\gdef\tn@single{#1}},
         tn index/.code={\tn@indexkey{#1}},
         tn@tname/.code={\gdef\tn@tname{#1}}}
% tn index: the list as \tn@index (<name>:<side>,...), the sides it names as
% \tn@legs, each once, and how many indices as \tn@rank.
\def\tn@indexkey#1{%
  \xdef\tn@ixall{\zap@space#1 \@empty}\gdef\tn@legs{}\gdef\tn@ixnames{}%
  \gdef\tn@index{}\gdef\tn@ixdirs{}\global\tn@ixn=0
  \foreach \tn@ii in \tn@ixall{\expandafter\tn@indexone\tn@ii::\tn@end}%
  \xdef\tn@rank{\the\tn@ixn}}
\newcount\tn@ixn
% an index <name>:<side>[:<way>], kept as <name>:<side> in \tn@index and
% <name>=<way> in \tn@ixdirs
\def\tn@indexone#1:#2:#3:#4\tn@end{%
  \global\advance\tn@ixn by 1
  \def\tn@ixs{#2}%
  \ifx\tn@index\@empty \xdef\tn@index{#1:#2}\else \xdef\tn@index{\tn@index,#1:#2}\fi
  \def\tn@ixw{#3}%
  \ifx\tn@ixw\@empty\else
    \tn@checkway{#3}{tn index: index `#1'}%
    \ifx\tn@ixdirs\@empty \xdef\tn@ixdirs{#1=#3}\else \xdef\tn@ixdirs{\tn@ixdirs,#1=#3}\fi
  \fi
  \ifx\tn@ixs\@empty
    \PackageError{tikz-tensors}{tn index: `#1' says no side; an index is
      <name>:<side>}{}\fi
  \edef\tn@ixt{\noexpand\in@{,#1,}{,\tn@ixnames,}}\tn@ixt
  \ifin@ \PackageError{tikz-tensors}{tn index: two indices are named `#1'}{}\fi
  \xdef\tn@ixnames{\tn@ixnames,#1}%
  \edef\tn@ixt{\noexpand\in@{,#2,}{,\tn@legs,}}\tn@ixt
  \ifin@\else \ifx\tn@legs\@empty \xdef\tn@legs{#2}\else \xdef\tn@legs{\tn@legs,#2}\fi\fi}
\def\tn@sin{in}\def\tn@sout{out}
\def\tn@checkway#1#2{\def\tn@cw{#1}%
  \ifx\tn@cw\tn@sin\else\ifx\tn@cw\tn@sout\else
    \PackageError{tikz-tensors}{#2 runs `#1'; an index runs in or out}{}\fi\fi}
\tn@renamed{/tikz/tn leg anchor}{tn leg anchor}{it is tn site anchor}
% Read a type's index properties without drawing it: \tn@legat, \tn@flux,
% \tn@single, \tn@index (and from it \tn@legs, \tn@rank and \tn@ixdirs)
% and \tn@tname afterwards hold what the style <#1> sets.
\def\tn@props#1{%
  \gdef\tn@legat{center}\gdef\tn@flux{}\gdef\tn@fluxw{}\gdef\tn@single{}%
  \gdef\tn@legs{}%
  \gdef\tn@rank{}\gdef\tn@tname{}\gdef\tn@index{}\gdef\tn@ixdirs{}%
  \begingroup\tikzset{#1}\endgroup}

% Each shape carries the size at which it reads as the same weight as the
% others -- a triangle has to be larger than a box to hold a label, a diamond a
% little less -- and paper fill, so that an index drawn to its center is covered
% by it. A fill of a convention's own replaces the paper. A triangle's index
% leaves from a corner of its flat side -- the index continues that side's own
% line rather than leaving from the middle of a slanted edge, which is how a
% canonical form is read; pointing right that is corner 3, pointing left corner
% 2 (one shape at a 180 degree turn, so no single anchor names both).
\tikzset{
  tn box/.style            = {tn node, rectangle, fill=grey1, minimum size=7mm,
                              inner sep=1pt},
  tn circle/.style         = {tn node, circle, fill=grey1, minimum size=8mm},
  tn capsule/.style        = {tn node, rectangle, tn@radius=4mm,
                              fill=grey1, minimum size=8mm},
  tn rounded/.style        = {tn node, rectangle, tn@radius=1.2mm,
                              fill=grey1, minimum size=7.5mm, inner sep=1pt},
  tn triangle right/.style = {tn node, regular polygon, regular polygon sides=3,
                              shape border rotate=-90, fill=grey1,
                              minimum size=10.5mm, font=\small,
                              tn site anchor=corner 3,
                              tn@lean={left=in, right=out}},
  tn triangle left/.style  = {tn triangle right, shape border rotate=90,
                              tn site anchor=corner 2,
                              tn@lean={right=in, left=out}},
  tn triangle up/.style    = {tn node, isosceles triangle,
                              isosceles triangle stretches, shape border rotate=90,
                              fill=grey1, minimum width=10.5mm, minimum height=9mm,
                              font=\small, tn single=up,
                              tn@lean={down=in, up=out}},
  tn triangle down/.style  = {tn triangle up, shape border rotate=-90,
                              tn single=down, tn@lean={up=in, down=out}},
  tn diamond/.style        = {tn node, diamond, fill=grey1, minimum size=9.5mm,
                              font=\small},
  tn wide/.style           = {minimum width=17mm},
  tn tall/.style           = {minimum height=17mm},
  tn small/.style          = {minimum size=6mm, font=\scriptsize},
  tn flat/.style           = {minimum width=10.5mm, minimum height=6mm},
  tn dot/.style            = {tn node, circle, fill=stroke, minimum size=1.8mm},
  tn oplus/.style          = {tn node, circle, fill=grey1, minimum size=5mm,
                              path picture={\draw[tn line, draw=stroke]
                                (path picture bounding box.north)
                                  -- (path picture bounding box.south)
                                (path picture bounding box.west)
                                  -- (path picture bounding box.east);}},
  tn frame/.style          = {tn line, draw=black!50, dashed, rounded corners=2mm},
}

% Two types are named for the structure that draws them rather than for a
% shape, because a layout needs something to draw with and a convention needs
% somewhere to colour it. Both are plain outlines here.
%   tn gate  a gate on indices that run through it (also a MERA's disentangler)
%   tn env   the block at either end of a stack: the rest of the network,
%            already contracted -- the L and R of an effective Hamiltonian, the
%            fixed points of a uniform state. It spans every layer of its stack
%            and has one index in each.
\tikzset{
  tn gate/.style = {tn rounded, inner sep=3pt, font=\small},
  tn env/.style  = {tn box, minimum width=\tn@envw, font=\small},
}

% The way each side of tensor <#1> runs, n@<t>@fx@<side>, from \tn@flux, and
% the way its shape leans, n@<t>@wfx@<side>, from \tn@fluxw (tn@lean); and
% each named index's own, n@<t>@ifx@<name>, from \tn@ixdirs.
\def\tn@setflux#1{%
  \foreach \tn@fs in {left,right,up,down,nw,ne,se,sw}{\tn@set{n@#1@fx@\tn@fs}{}%
    \tn@set{n@#1@wfx@\tn@fs}{}}%
  \foreach \tn@fp in \tn@fluxw{\expandafter\tn@fluxone\tn@fp=\tn@end{#1}{wfx}}%
  \foreach \tn@fp in \tn@flux{\expandafter\tn@fluxone\tn@fp=\tn@end{#1}{fx}}%
  \foreach \tn@fp in \tn@ixdirs{\expandafter\tn@fluxone\tn@fp=\tn@end{#1}{ifx}}}
\def\tn@fluxone#1=#2=#3\tn@end#4#5{%
  \tn@checkway{#2}{tn flux: side `#1'}\tn@set{n@#4@#5@#1}{#2}}

% ---- placing a node ----------------------------------------------------------
% \tn@place{<type>}{<name>}{<x>}{<y>}{<label>}
% How every layout puts a tensor down: at (<x>,<y>), or, when its type gives
% its index an anchor of its own, moved across so that the anchor is there --
% a probe of the same type, drawn nowhere, measures how far. The way its sides
% run and its single side are kept with the node (n@<name>@fx@<side>,
% n@<name>@one), for the routes that later meet it.
\def\tn@place#1#2#3#4#5{%
  \tn@props{#1}%
  \tn@set{n@#2@one}{\tn@single}\tn@set{n@#2@legs}{\tn@legs}%
  \tn@setflux{#2}%
  \tn@set{n@#2@type}{\detokenize{#1}}\tn@set{n@#2@rank}{\tn@rank}%
  \tn@set{n@#2@tname}{\tn@tname}\tn@set{n@#2@index}{\tn@index}%
  \ifx\tn@legat\tn@center
    \node[#1] (#2) at (#3,#4) {#5};\tn@tracenode{#2}%
  \else
    \path (0,0) node[#1, overlay, draw=none, fill=none] (tn@probe) {};
    \tn@anc{tn@probe}{\tn@legat}\tn@d=\tn@ax\relax
    \node[#1] (#2) at ($(#3,#4)-(\the\tn@d,0)$) {#5};\tn@tracenode{#2}%
  \fi}
